Les chaînes logistiques modernes opérant dans les réseaux de l'Internet Physique font face à un écart fondamental de qualité de signal~: l'urgence déclarée par les opérateurs diverge souvent de l'urgence réelle physiquement calculable. Il en résulte du sur-triage (gaspillage de ressources) et du sous-triage (risque de rupture caché). Dans des contextes Supply Chain~4.0 marqués par la volatilité et les dépendances multi-rangs, aucune architecture de jumeau numérique existante ne corrige pleinement ce biais cognitif tout en le propageant de manière cohérente sur un réseau de fournisseurs.

Ce mémoire synthétise un stage de R\&D qui propose un score d'adéquation cognitive $A(t) \in [0, 100]$ pour mesurer et corriger l'écart entre l'urgence déclarée par l'opérateur et une urgence réelle $U_r(t)$ multi-blocs, calculée par le jumeau numérique et propagée sur le graphe orienté acyclique du réseau. L'approche combine une élicitation par comparaisons par paires (AHP, Fuzzy Best-Worst Method) pour l'urgence déclarée, une agrégation noisy-OR de six blocs de KPI opérationnels pour l'urgence réelle, le surclassement PROMETHEE~II et une simulation systématique de chocs pour la criticité réseau, ainsi qu'un moteur d'événements calibrés pour maintenir à jour l'état des KPI sous-jacents. Le cadre est implémenté en Python et opéré comme un jumeau numérique persistant et multi-tenant, exercé à travers un serious game à cadence hebdomadaire, avec un service de calibration dédié comparant le signal d'adéquation aux résultats opérationnels enregistrés.
